% vim: ai sw=2 spell spelllang=cs

\begin{frame}
  \frametitle{Komunikace s HW}

  \begin{itemize}
    \item Obvykle práce ovladačů v jádře

    \pause

    \item Ale často není jádro třeba
    \begin{itemize}
      \item Zařízení na sériovém portu
      \item Skenery, tiskárny a podobná USB zařízení
      \item I jednoduchá PCI zařízení
    \end{itemize}
  \end{itemize}

\end{frame}

\begin{frame}
  \frametitle{Různé přístupy ke komunikaci s HW}

  \begin{enumerate}
    \item Jádro HW ovládá a uživatelovi nevystavuje
    \begin{itemize}
      \item V Linuxu téměř nic
    \end{itemize}

    \pause

    \item Jádro zaštiťuje nízkoúrovňovou komunikaci
    \begin{itemize}
      \item Zařízení v \file{/dev/} (\file{ttyS*}, \file{ttyUSB*}, \file{lp*})
      \item Soubory v \file{/sys/} (gadgety)
    \end{itemize}

    \pause

    \item Jádro obsluhuje jen přerušení
    \begin{itemize}
      \item Vrstva UIO
      \item Dokumentace: \href{http://blackfin.uclinux.org/doku.php?id=auto_generated_kernel_docs:uio-howto}{uio-howto}
    \end{itemize}

    \pause

    \item Jádro nedělá s HW nic
    \begin{itemize}
      \item Uživatel používá instrukce/volání
      \begin{itemize}
	\item Získání povolení: \code{ioperm} nebo \code{iopl}
	\item Přímý přístup: \code{out+in}
      \end{itemize}

      \item Soubory v \file{/sys/} (PCI)
      \begin{itemize}
	\item \code{mmap}
      \end{itemize}

      \item Knihovny na přímou komunikaci s HW
      \begin{itemize}
	\item \file{libusb}, \file{libpciaccess}
      \end{itemize}
    \end{itemize}
  \end{enumerate}

\end{frame}

\mysection{I/O porty}

\begin{frame}{Přístup instrukcemi/voláními}{I/O porty}

  \begin{itemize}
    \item Porty jsou specifikum některých CPU (včetně \xes)
    \begin{itemize}
      \item Samostatná (malá) sběrnice
      \item Speciální instrukce (\code{in}, \code{out} na \xes)
      \item Je zde řadič klávesnice, PC spkr, časovače a ovladače přerušení,
	ladicí port (\code{0x80} $\Rightarrow$ segmentový displej na desce),
	\dots
    \end{itemize}

    \pause

    \item Žádost o přístup: \code{ioperm}, \code{iopl}
    \item Čtení: \code{inX}
    \item Zápis: \code{outX}, kde X $\in\left\{ \text{\code{b}},\,
      \text{\code{w}},\, \text{\code{l}} \right\}$ 
    \item Dokumentace: \cmd{man ioperm} a \cmd{man inb} 
  \end{itemize}

\end{frame}

\begin{frame}
  \frametitle{Úkol}

  \heading{Zahrajte melodii na PC speaker} (doplňte
    \file{pb173-bin/11/speaker.c})
  \begin{enumerate}
    \item Nastartujte si virtuální stroj (\file{pb173-bin/qemu-start})
    \item Dopište tělo funkce \code{beep}
    \begin{itemize}
      \item Zapište \code{0xb6} na port \code{0x43}
      \item Zapište \code{div} na port \code{0x42}
      \item Zapište \code{(div >> 8)} na port \code{0x42}
      \item Přečtěte port \code{0x61} (obsah spodních 2 bitů zahoďte)
      \item Zapište přečtenou hodnotu \code{| 0x03} zpět na port \code{0x61}
    \end{itemize}

    \item Dopište tělo funkce \code{stop_beep}
    \begin{itemize}
      \item Přečtěte port \code{0x61}
      \item Zapište \code{0} na spodní 2 bity portu \code{0x61}
    \end{itemize}

    \item Po všech zápisech a čteních čekejte alespoň 10 $\mu$s
    \begin{itemize}
      \item Máte k dispozici \code{usleep}
    \end{itemize}

    \item Přidejte do \code{main} před \code{play} volání \code{ioperm}
    \item Spusťte
  \end{enumerate}

\end{frame}

\mysection{PCI zařízení}

../../pb173/08/pci.tex

\mysection{\code{mmap} pro přímou komunikaci}

\begin{frame}
  \frametitle{\code{mmap} pro přímou komunikaci s PCI}

  \begin{itemize}
    \item Každé PCI zařízení může mít:
    \begin{itemize}
      \item I/O porty (komunikace viz dříve)
      \item \emph{Paměť ve fyzickém prostoru}
      \item Přerušení, DMA, atd. (UIO)
    \end{itemize}

    \pause
    \medskip

    \item Soubory \file{/sys/bus/pci/devices/*/resource*}
    \item Mapují se pomocí \code{mmap}
    \item Obsah: ve specifikaci konkrétního zařízení
  \end{itemize}

\end{frame}

\begin{frame}
  \frametitle{Úkol}

  \heading{Komunikace s EDU zařízením}
  \begin{enumerate}
    \item Zjistěte číslo EDU karty z výstupu \cmd{lspci}
    \begin{itemize}
      \item Má ID \code{0x11e8} v \cmd{qemu}
      \item Hledáte něco jako \code{01:01.0}
    \end{itemize}

    \item Mapujte \file{resource0} EDU zařízení
    \begin{itemize}
      \item V \file{/sys/bus/pci/devices/...}
      \item Jen v qemu s \cmd{\-device edu}
      \item Mapujte jako \code{volatile uint32_t *}
    \end{itemize}

    \item Vypište, co je na první a druhé pozici
    \item Zapište nějaké číslo na druhou pozici
    \item Vypište, co je na druhé pozici
    \item Spusťte
  \end{enumerate}

\end{frame}

\mysection{Knihovny pro přímou komunikaci s HW}

\begin{frame}
  \frametitle{Knihovny pro přímou komunikaci s HW}

  \begin{itemize}
    \item Kromě přímého \code{mmap} přístupu existují i knihovny
    \item \file{libusb}, \file{libpciaccess}
    \begin{itemize}
      \item Jsou přenositelné
      \item Poskytují lepší rozhraní
      \item Pracují se soubory v \file{/proc} a \file{/sys}
    \end{itemize}
  \end{itemize}

\end{frame}

\begin{frame}[fragile]{\file{libpciaccess}}

  \begin{itemize}
    \item \cmd{gcc ... -lpciaccess} (\header{pciaccess.h})

    \pause

    \item Inicializace: \code{pci_system_init}
    \item Deinicializace: \code{pci_system_cleanup}

    \pause
    \medskip

    \item Vytvoření iterátoru (\code{struct pci_device_iterator *})
    \begin{itemize}
      \item Vím ID: \code{pci_id_match_iterator_create} a
	\code{struct pci_id_match}
      \item Vím slot: \code{pci_slot_match_iterator_create} a
	\code{struct pci_slot_match}
    \end{itemize}

    \pause

    \item Procházení iterátoru: \code{pci_device_next}
    \begin{itemize}
      \item Vrací \code{struct pci_device *} (a \code{NULL} nakonec)
    \end{itemize}

    \pause

    \item Zničení iterátoru: \code{pci_iterator_destroy}
  \end{itemize}

  \onslide<3->

  \begin{block}{Struktury \code{pci_id_match} a \code{pci_slot_match}}
    \begin{tabular}{lp{6em}lp{6em}l}
      \begin{lstlisting}
uint32_t vendor_id;
uint32_t device_id;
uint32_t subvendor_id;
uint32_t subdevice_id;
      \end{lstlisting} &&
      \begin{lstlisting}
uint32_t domain;
uint32_t bus;
uint32_t dev;
uint32_t func;
      \end{lstlisting} &&
      \begin{lstlisting}
PCI_MATCH_ANY
      \end{lstlisting}
    \end{tabular}
  \end{block}


\end{frame}

\begin{frame}
  \frametitle{Úkol}

  \heading{Výpis PCI zařízení pomocí \file{libpciaccess}}
  \begin{enumerate}
    \item Volejte \code{pci_system_init}
    \item Vytvořte iterátor
    \begin{itemize}
      \item S parametrem \code{NULL}, tj.\ projdi všechny
    \end{itemize}

    \item Projděte zařízení

    \item Pro každé vypište číslo
    \begin{itemize}
      \item Domény
      \item Sběrnice
      \item Zařízení
      \item Funkce
    \end{itemize}

    \item Spusťte a zkontrolujte s \cmd{lspci -nn}
  \end{enumerate}

\end{frame}

\begin{frame}
  \frametitle{\file{libpciaccess} a zařízení}

  \begin{itemize}
    \item Před další prací je nutné inicializovat \code{struct pci_device}
    \begin{itemize}
      \item Pomocí \code{pci_device_probe}
      \item Potom lze pracovat s dalšími prvky \code{struct pci_device}
      \begin{itemize}
        \item \code{vendor_id}, \code{device_id}, \code{irq}, \dots
      \end{itemize}
    \end{itemize}

    \pause
    \medskip

    \item Mapování regionu: \code{pci_device_map_range}
    \begin{itemize}
      \item Mapuje region do paměti
      \item Podobné \code{mmap} regionu z dnešního cvičení
    \end{itemize}

    \pause
    \medskip

    \item Konfigurační prostor
    \begin{itemize}
      \item Zapisuje BIOS a OS
      \item Informace o zařízení
      \item Čtení: \code{pci_device_cfg_read*}
      \item Zápis: \code{pci_device_cfg_write*}
    \end{itemize}
  \end{itemize}

\end{frame}

\begin{frame}
  \frametitle{Úkol}

  \heading{Práce s EDU pomocí \file{libpciaccess}}
  \begin{enumerate}
    \item Mapujte si 0.\ region
    \begin{itemize}
      \item \code{pci_device_map_range}
      \item Mapujte jako \code{volatile uint32_t *}
      \item Podobně jako \code{mmap} předtím
    \end{itemize}

    \item Vypište, co je na první a \emph{třetí} pozici
    \item Zapište nějaké malé číslo na třetí pozici
    \item Počkejte, než bude osmá pozice \code{0}
    \item Vypište, co je na třetí pozici
    \item Spusťte
  \end{enumerate}

\end{frame}
